\section{Proof of Theorem~\ref{thm:common-invariant-limit}}
\label{sec:convergence}

% The previous version of this section is preserved verbatim in
% paper/sections/convergence.tex while the section is rewritten.

The proof starts from the patch representation. For~$A\Subset\Lambda$ and law
$\mu$ we have
\begin{equation*}
(\mu P_t)(\chi_A)
=
\E_A\Cb{
\prod_{P\in\mathcal B_t}C(P)\,
\mu\bb{
\prod_{P\in\mathcal E_t}C\bb{\eta\bb{i(P)},P}
}
},
\end{equation*}
where the variables~$\eta(i(P))$ come from the initial configuration in the
backward dual representation. We aim to show this converges to
\begin{equation*}
\E_A\Cb{
\prod_{P\in\mathcal P}C(P)
\ind\bb{\abs{\mathcal P}<\infty}
}.
\end{equation*}
The event~$\cb{\abs{\mathcal P}<\infty}$ is the event that the dual process eventually dies, and from that point there are no more successful interactions. In particular, on this event the
limiting product contains no end-patch contribution and hence no dependence on
the initial configuration. This suggests proving convergence by showing first
that the contribution from the event that a successful interaction occurs far
in time is negligible.

The comparison with pure deaths introduced earlier gives exactly this estimate.
Section~\ref{sec:signed-dual} showed that adding pure deaths at rate
$\varepsilon$ leaves the dual jump process unchanged and decreases its
Feynman--Kac potential by~$\varepsilon\abs{A_s}$. Corollary~\ref{cor:pure-death-comparison}
then used the same observation at the level of patches: removing these deaths
leaves the successful interactions and patch families unchanged and increases
the patch contributions. If a successful interaction occurs at time~$u$, its
ancestry consists of active patches whose lifetimes add up to~$u$. The
contribution for the original system is therefore smaller, relative to the system
with the pure deaths removed, by a factor~$e^{-\varepsilon u}$. Since the latter
provides an integrable upper bound, the contribution from the event that a
successful interaction occurs after~$T$ is of order~$e^{-\varepsilon T}$.

It remains to consider the event that there is no successful interaction after
$T$ and show that on this event the dependence on initial law decays. The end patches present at a later horizon~$t$ still contain the factors
$C(z,P)$ through which the initial configuration enters the representation.
Once there are no further successful interactions, the dependence of each such
factor on~$z$ decays at rate at least~$\varepsilon$, and hence is of order
$e^{-\varepsilon(t-T)}$. This is the mechanism by which the dependence on the
initial law disappears. The resulting estimate, however, gets worse with
the number of end patches.

For this reason we also control the event that there are many end patches. If
the interaction cone up to time~$T$ stays in a finite region~$R$, there are at
most~$\abs R$ end patches. On the complementary event the interaction cone exits~$R$ before time~$T$, and its contribution is controlled by finite propagation. We take~$R$ with radius
proportional to~$T$. Polynomial growth ensures~$\abs R\lesssim(1+T)^D$, while the
finite-propagation error is exponentially small. In the proof below we make these
three estimates precise and then combine them with~$T=t/2$. We first prove the convergence estimate for initial laws in~$\mathcal M_*$ and then extend it to~$\mathcal M_-$ by writing each measure in~$\mathcal M_{-,K}$ as an affine combination of two measures in~$\mathcal M_*$.

\subsection{Three estimates}
\label{subsec:convergence-three-estimates}

The estimates below use the nonnegativity of the patch contributions and are
therefore proved for initial laws in~$\mathcal M_*$. This is enough for the
larger class in Theorem~\ref{thm:common-invariant-limit}. Indeed, for
$\mu\in\mathcal M_{-,K}$, set
\begin{equation*}
\overline\mu
=
\frac{\mu+K\mu_{\mb 1}}{1+K}
\in\mathcal M_*.
\end{equation*}
Then~$\mu=(1+K)\overline\mu-K\mu_{\mb 1}$, and~$\mu_{\mb 1}\in\mathcal M_*$. Hence any bound
\begin{equation*}
\sup_{\nu\in\mathcal M_*}
\abs{(\nu P_t)(f)-\ell(f)}
\leq B_f(t)
\end{equation*}
implies
\begin{equation}
\label{eq:m-minus-reduction}
\sup_{\mu\in\mathcal M_{-,K}}
\abs{(\mu P_t)(f)-\ell(f)}
\leq
(1+2K)B_f(t).
\end{equation}
It therefore suffices to prove the convergence estimate uniformly on
$\mathcal M_*$.

For any probability measure~$\mu$, set
\begin{equation}
\label{eq:finite-horizon-weight}
W_t^\mu
=
\prod_{P\in\mathcal B_t}C(P)\,
\mu\bb{
 \prod_{P\in\mathcal E_t}
 C\bb{\eta\bb{i(P)},P}
},
\end{equation}
and set
\begin{equation*}
W
=
\prod_{P\in\mathcal P}C(P)
\ind\bb{\abs{\mathcal P}<\infty}.
\end{equation*}

For~$0\leq T<t$, let
\begin{align}
\label{eq:no-late-interaction-events}
L_{T,t}
&=
\cb{\text{no successful interaction occurs in }(T,t]},\notag\\
L_T
&=
\cb{\text{no successful interaction occurs after }T}.
\end{align}
The first estimate controls the contribution of event with late successful interactions.
\begin{lemma}[Late interactions]
\label{lem:late-interactions}
For~$\mu\in\mathcal M_*$ and~$0\leq T<t$,
\begin{align}
0
&\leq
\E_A\Cb{W_t^\mu\ind\bb{L_{T,t}^c}}
\leq
e^{-\varepsilon T},
\label{eq:finite-late-interaction-bound}
\\
0
&\leq
\E_A\Cb{W\ind\bb{L_T^c}}
\leq
e^{-\varepsilon T}.
\label{eq:limiting-late-interaction-bound}
\end{align}
\end{lemma}

We next restrict to the event that there is no successful interaction after~$T$.
The error from replacing an end-patch contribution at time~$t$ by its limit
accumulates over the end patches. To bound their number, let
\begin{equation*}
\mb{Cone}_T
=
\bigcup_{(i,u,S)\in\mathcal I_T}
\bb{\cb{i}\cup S}\setminus\cb{\infty}
\end{equation*}
be the set of sites appearing in the successful-interaction skeleton up to
time~$T$, and for~$A\subseteq R\Subset\Lambda$ set
\begin{equation}
\label{eq:confinement-event}
E_T^R=\cb{\mb{Cone}_T\subseteq R}.
\end{equation}
On~$E_T^R$ there are at most~$\abs R$ end patches at time~$T$.
\begin{lemma}[Relaxation after the last interaction]
\label{lem:no-late-relaxation}
For~$A\subseteq R\Subset\Lambda$,~$\mu\in\mathcal M_*$, and~$0\leq T<t$,
\begin{equation}
\label{eq:no-late-relaxation-bound}
\abs{
 \E_A\Cb{W_t^\mu\ind\bb{E_T^R\cap L_{T,t}}}
 -
 \E_A\Cb{W\ind\bb{E_T^R\cap L_T}}
}
\leq
\bb{1+\abs R}e^{-\varepsilon\bb{t-T}}.
\end{equation}
\end{lemma}

It remains to control the event~$(E_T^R)^c$. On~$E_T^R$, the dual during the
first~$T$ units of time stays inside~$R$, so its contribution is unchanged
if the spins outside~$R$ are fixed at~$0$. Let~$P_t^{R,0}$ be the semigroup of
this zero-boundary dynamics on~$R$, and set
\begin{equation}
\label{eq:confinement-error}
\rho_A(T,R)
=
\norm{\bb{P_T-P_T^{R,0}}\chi_A}_\infty.
\end{equation}
\begin{lemma}[Spatial confinement]
\label{lem:spatial-confinement}
For~$A\subseteq R\Subset\Lambda$,~$\mu\in\mathcal M_*$, and~$0\leq T<t$,
\begin{align}
0
&\leq
\E_A\Cb{W_t^\mu\ind\bb{(E_T^R)^c}}
\leq
\rho_A(T,R),
\label{eq:finite-horizon-confinement-bound}
\\
0
&\leq
\E_A\Cb{W\ind\bb{(E_T^R)^c}}
\leq
\rho_A(T,R).
\label{eq:limiting-confinement-bound}
\end{align}
\end{lemma}

\subsection{Completion of the proof}
\label{subsec:convergence-completion}

\begin{proof}[Proof of Theorem~\ref{thm:common-invariant-limit}]
Fix~$A\Subset\Lambda$,~$\mu\in\mathcal M_*$,~$0\leq T<t$, and a finite
$R\supseteq A$. Since
\begin{equation*}
1
=
\ind\bb{L_{T,t}^c}
+
\ind\bb{E_T^R\cap L_{T,t}}
+
\ind\bb{(E_T^R)^c\cap L_{T,t}},
\end{equation*}
and the analogous identity holds with~$L_T$ in place of~$L_{T,t}$, the patch
representation gives
\begin{align*}
(\mu P_t)(\chi_A)-\E_A\Cb{W}
&=
\E_A\Cb{W_t^\mu\ind\bb{L_{T,t}^c}}
-
\E_A\Cb{W\ind\bb{L_T^c}}
\\
&\quad+
\E_A\Cb{W_t^\mu\ind\bb{E_T^R\cap L_{T,t}}}
-
\E_A\Cb{W\ind\bb{E_T^R\cap L_T}}
\\
&\quad+
\E_A\Cb{W_t^\mu\ind\bb{(E_T^R)^c\cap L_{T,t}}}
-
\E_A\Cb{W\ind\bb{(E_T^R)^c\cap L_T}}.
\end{align*}
Since~$W_t^\mu,W\geq0$ and
$(E_T^R)^c\cap L_{T,t}\subseteq(E_T^R)^c$, the three lemmas above give
\begin{equation}
\label{eq:three-term-convergence-bound}
\abs{(\mu P_t)(\chi_A)-\E_A\Cb{W}}
\leq
2e^{-\varepsilon T}
+
\bb{1+\abs R}e^{-\varepsilon\bb{t-T}}
+
2\rho_A(T,R).
\end{equation}

Apply Lemma~\ref{lem:finite-propagation} with~$a=\varepsilon$. There are
$v<\infty$ and~$C_A<\infty$ such that, for
\begin{equation*}
R_T=B\bb{A,\lceil vT\rceil},
\end{equation*}
we have
\begin{equation*}
\rho_A(T,R_T)\leq C_Ae^{-\varepsilon T}.
\end{equation*}
Polynomial growth gives, after increasing~$C_A$ if necessary,
\begin{equation*}
\abs{R_T}\leq C_A(1+T)^D.
\end{equation*}
Taking~$T=t/2$ in~\eqref{eq:three-term-convergence-bound} therefore yields
\begin{equation}
\label{eq:m-star-monomial-convergence}
\sup_{\mu\in\mathcal M_*}
\abs{
(\mu P_t)(\chi_A)
-
\E_A\Cb{W}
}
\leq
K_A(1+t)^D e^{-\varepsilon t/2}
\end{equation}
for some~$K_A<\infty$.

We now identify these limits with a probability measure. Since
$\mu_{\mb 1}\in\mathcal M_*$ and~$\Omega$ is compact,
$\mu_{\mb 1}P_t$ has a weakly convergent subsequence. Let~$\pi$ be one such
subsequential limit. Equation~\eqref{eq:m-star-monomial-convergence} gives
\begin{equation}
\label{eq:pi-patch-moments}
\pi(\chi_A)
=
\E_A\Cb{
\prod_{P\in\mathcal P}C(P)
\ind\bb{\abs{\mathcal P}<\infty}
}
\qquad
\text{for every }A\Subset\Lambda.
\end{equation}
The monomials determine every local function, so every subsequential limit
has the same finite-dimensional marginals. Hence
$\mu_{\mb 1}P_t\Rightarrow\pi$ as~$t\to\infty$.

The measure~$\pi$ is invariant. Indeed, if~$g$ is local and~$s\geq0$, the
Feller property gives
\begin{equation*}
(\pi P_s)(g)
=
\lim_{t\to\infty}(\mu_{\mb 1}P_t)(P_sg)
=
\lim_{t\to\infty}(\mu_{\mb 1}P_{t+s})(g)
=
\pi(g).
\end{equation*}
Thus~$\pi P_s$ and~$\pi$ agree on all local functions and therefore are equal.

Every local function is a finite linear combination of monomials. Hence
\eqref{eq:m-star-monomial-convergence} and~\eqref{eq:pi-patch-moments} give, for
every local~$f$,
\begin{equation}
\label{eq:m-star-local-convergence}
\sup_{\mu\in\mathcal M_*}
\abs{(\mu P_t)(f)-\pi(f)}
\leq
K_f(1+t)^D e^{-\varepsilon t/2}
\end{equation}
for some~$K_f<\infty$. Applying~\eqref{eq:m-minus-reduction} with
$B_f(t)=K_f(1+t)^D e^{-\varepsilon t/2}$ proves
\eqref{eq:uniform-m-minus-convergence}. Since
$\mathcal M_-=\bigcup_{K\geq0}\mathcal M_{-,K}$, it also gives
$\mu P_t\Rightarrow\pi$ for every~$\mu\in\mathcal M_-$. Finally,
\eqref{eq:pi-patch-moments} is exactly
\eqref{eq:full-patch-limit-moments}.
\end{proof}

We also verify the claim following Theorem~\ref{thm:common-invariant-limit}
that every configuration with finitely many zeros gives an initial law in
$\mathcal M_-$. The pure-death assumption gives
\begin{equation*}
\lambda_i(\mb 0)\bb{1-\rho_i(\mb 0)}\geq\varepsilon,
\end{equation*}
so~$\rho_i(\mb 0)<1$. Since~$p_i^\star\leq\rho_i(\mb 0)$, we have
$p_i^\star<1$ for every~$i$. Let~$\eta$ have finite zero set
\begin{equation*}
F=\cb{i:\eta(i)=0}.
\end{equation*}
For every finite~$A$,
\begin{equation*}
\frac{
\abs{\delta_\eta(\chi_A^*)}
}{
\mu_{\mb 1}(\chi_A^*)
}
=
\prod_{i\in A\cap F}
\frac{p_i^\star}{1-p_i^\star}.
\end{equation*}
Consequently,
\begin{equation*}
K
\geq
\prod_{i\in F}
\max\cb{1,\frac{p_i^\star}{1-p_i^\star}}
\end{equation*}
ensures
\begin{equation*}
\delta_\eta(\chi_A^*)
+K\mu_{\mb 1}(\chi_A^*)
\geq0
\qquad
\text{for every }A\Subset\Lambda.
\end{equation*}
Thus~$\delta_\eta\in\mathcal M_{-,K}\subseteq\mathcal M_-$.

\begin{proof}[Proof of Corollary~\ref{cor:uniform-ergodicity}]
Suppose~$\mb p^\star\leq\mbs{\tfrac12}$. For every configuration~$\eta$ and
finite~$A$,
\begin{equation*}
\chi_A^*(\eta)
\geq
-\prod_{i\in A}\bb{1-p_i^\star}
=
-\mu_{\mb 1}(\chi_A^*),
\end{equation*}
because~$p_i^\star\leq1-p_i^\star$ at every site. Hence, for every probability
measure~$\mu$,
\begin{equation*}
\frac{\mu+\mu_{\mb 1}}2\bb{\chi_A^*}\geq0
\qquad
\text{for every }A\Subset\Lambda,
\end{equation*}
so~$\mu\in\mathcal M_{-,1}$. Theorem~\ref{thm:common-invariant-limit} therefore
gives the uniform bound~\eqref{eq:uniform-pointwise-convergence}, after
absorbing the fixed factor~$3$ into~$K_f$. In particular, every initial law
converges to~$\pi$, so~$\pi$ is the unique invariant measure.
\end{proof}

\subsection{Proofs of the estimates}
\label{subsec:convergence-estimate-proofs}

We now prove the three estimates used above.

\begin{proof}[Proof of Lemma~\ref{lem:late-interactions}]
Write
\begin{equation}
\label{eq:uniform-death-decomposition}
\Lcal
=
\Lcal^\varepsilon
+
\varepsilon\Ncal^{\mb 0},
\end{equation}
and let~$P_s^\varepsilon$ and~$C^\varepsilon$ denote the semigroup and patch
contributions for~$\Lcal^\varepsilon$. In the proof of
Corollary~\ref{cor:pure-death-comparison}, we show that the system with pure deaths removed has the same dual process, successful-interaction skeleton, and patch families. For a
probability measure~$\mu$, define
\begin{align*}
W_t^{\varepsilon,\mu}
&=
\prod_{P\in\mathcal B_t}C^\varepsilon(P)\,
\mu\bb{
 \prod_{P\in\mathcal E_t}
 C^\varepsilon\bb{\eta\bb{i(P)},P}},\\
W^\varepsilon
&=
\prod_{P\in\mathcal P}C^\varepsilon(P)
\ind\bb{\abs{\mathcal P}<\infty}.
\end{align*}
For every patch and every~$\mu\in\mathcal M_*$, respectively,
\begin{equation*}
0\leq C(P)\leq C^\varepsilon(P)
\end{equation*}
and
\begin{equation*}
0\leq
\mu\bb{
 \prod_{P\in\mathcal E_t}C\bb{\eta\bb{i(P)},P}
}
\leq
\mu\bb{
 \prod_{P\in\mathcal E_t}C^\varepsilon\bb{\eta\bb{i(P)},P}
}.
\end{equation*}
These inequalities are proved in the proof of
Corollary~\ref{cor:pure-death-comparison}. Consequently,
\begin{equation}
\label{eq:death-removed-weight-comparison}
0\leq W_t^\mu\leq W_t^{\varepsilon,\mu},
\qquad
\E_A\Cb{W_t^{\varepsilon,\mu}}
=(\mu P_t^\varepsilon)(\chi_A)
\leq1.
\end{equation}

If a patch~$P$ is of type~$\mathsf{XO}$, meaning it ends with an outgoing successful interaction, then its site is
dual-active throughout~$[s(P),e(P))$. Removing the pure deaths raises the
one-site potential from~$V_i$ to~$V_i+\varepsilon$ and leaves the dual jump
process and the conditional patch law unchanged. Therefore
\begin{equation}
\label{eq:outgoing-terminal-exact-factor}
C(P)
=
e^{-\varepsilon\bb{e(P)-s(P)}}C^\varepsilon(P).
\end{equation}

We now show that for a successful interaction happening at time~$u>0$ there exists a trail of~$\mathsf{XO}$ type patches that spans time from~$0$ to~$u$. Every successful interaction ends a~$\mathsf{XO}$ patch. We call the interaction record that starts this patch the ancestor of the previous successful interaction. Jumping from ancestor to ancestor generates a finite sequence of interaction records ending at the formal initial interaction~$\iota_0$. Indeed, the boundary times decrease
strictly, and local finiteness of the graphical construction forces the chain
to reach time~$0$ after finitely many steps, call that number~$n$. This ancestry trail gives the sequence of patches
$P_1,\ldots,P_n$, each joining a successful interaction with its ancestor. All these patches are of type~$\mathsf{XO}$ and
\begin{equation}
\label{eq:outgoing-chain-total-length}
\sum_{k=1}^n\bb{e(P_k)-s(P_k)}=u.
\end{equation}

\begin{figure}[htbp]
\centering
% Schematic for the backward chain in Lemma 6.3.
% The highlighted vertical intervals are XO patches; the double-headed links
% are precisely the successful interactions used by the chain.
\begingroup
\definecolor{convtrailred}{RGB}{175,32,38}
\definecolor{convtrailoutline}{gray}{0.18}
\definecolor{convtrailcone}{gray}{0.88}
\definecolor{convtrailconeborder}{gray}{0.25}
\definecolor{convtrailsecondary}{gray}{0.43}
\begin{tikzpicture}[x=0.56cm,y=0.56cm]
\path[use as bounding box] (0.20,-0.12) rectangle (7.22,7.48);
\tikzset{
  convtrail cone/.style={fill=convtrailcone,draw=convtrailconeborder,
    line width=0.62pt,line join=round},
  convtrail patch/.style={draw=convtrailred,line width=2.65pt,
    line cap=round},
  convtrail interaction/.style={draw=convtrailoutline,line width=0.78pt,
    shorten <=1.2pt,shorten >=1.2pt,
    {Stealth[length=1.35mm,width=0.96mm]}-{Stealth[length=1.35mm,width=0.96mm]}},
  convtrail axis/.style={draw=convtrailoutline,line width=0.62pt,
    -{Stealth[length=1.8mm,width=1.2mm]}},
  convtrail tick/.style={draw=convtrailoutline,line width=0.52pt},
  convtrail baseline/.style={draw=convtrailoutline,line width=0.52pt},
  convtrail guide/.style={draw=convtrailsecondary,line width=0.48pt,
    dash pattern=on 2.2pt off 2.0pt},
  convtrail label/.style={font=\scriptsize,text=convtrailoutline,inner sep=1pt},
  convtrail secondary label/.style={font=\scriptsize,
    text=convtrailsecondary,inner sep=1pt}
}

% A deliberately irregular spacetime outline of the interaction cone.
\path[convtrail cone]
  (3.26,0)
  .. controls (3.07,0.55) and (2.95,1.08) .. (2.98,1.56)
  .. controls (2.99,1.91) and (2.45,2.17) .. (2.28,2.63)
  .. controls (2.10,3.11) and (2.48,3.48) .. (2.25,3.87)
  .. controls (2.00,4.31) and (1.45,4.50) .. (1.30,5.03)
  .. controls (1.15,5.55) and (0.82,6.01) .. (0.66,6.50)
  -- (5.56,6.50)
  .. controls (5.44,6.04) and (5.06,5.68) .. (5.25,5.18)
  .. controls (5.43,4.70) and (5.06,4.43) .. (4.87,4.08)
  .. controls (4.64,3.65) and (5.07,3.31) .. (4.92,2.87)
  .. controls (4.77,2.40) and (4.34,2.04) .. (4.47,1.58)
  .. controls (4.60,1.11) and (4.12,0.54) .. (3.95,0)
  -- cycle;

\node[convtrail secondary label,anchor=west] at (1.38,5.62)
  {$\mb{Cone}$};

% The backward chain.  Its patch lifetimes partition [0,u).
\draw[convtrail baseline] (0.40,0)--(5.78,0);
\draw[convtrail patch] (3.35,0)--(3.35,1.34);
\draw[convtrail interaction] (3.35,1.34)--(4.45,1.34);
\draw[convtrail patch] (4.45,1.34)--(4.45,2.66);
\draw[convtrail interaction] (2.75,2.66)--(4.45,2.66);
\draw[convtrail patch] (2.75,2.66)--(2.75,3.98);
\draw[convtrail interaction] (2.75,3.98)--(4.55,3.98);
\draw[convtrail patch] (4.55,3.98)--(4.55,5.18);
\draw[convtrail interaction] (3.25,5.18)--(4.55,5.18);
\draw[convtrail patch] (3.25,5.18)--(3.25,6.50);

% The prescribed successful interaction at time u.
\draw[convtrail interaction] (3.25,6.50)--(5.12,6.50);
\node[convtrail label,anchor=south] at (4.18,6.66) {$(i,u,S)$};

% Time axis and the guide locating the final interaction time.
\draw[convtrail guide] (5.12,6.50)--(6.38,6.50);
\draw[convtrail axis] (6.38,0)--(6.38,7.28);
\draw[convtrail tick] (6.28,0)--(6.48,0);
\draw[convtrail tick] (6.28,6.50)--(6.48,6.50);
\node[convtrail label,anchor=west] at (6.54,0) {$0$};
\node[convtrail label,anchor=west] at (6.54,6.50) {$u$};
\node[convtrail secondary label,anchor=south] at (6.38,7.29)
  {time};
\end{tikzpicture}
\endgroup

\caption{The ancestry of a successful interaction at time~$u$. Each dark
rectangle is a patch ending at an outgoing successful interaction. The patch lifetimes sum to~$u$.}
\label{fig:backward-outgoing-chain}
\end{figure}
\FloatBarrier

On~$L_{T,t}^c$, let~$u\in(T,t]$ be the time of the first successful interaction after~$T$
and let~$P_1,\ldots,P_n$ be its ancestry chain. Every patch in the chain belongs
to~$\mathcal B_t$, so~\eqref{eq:outgoing-terminal-exact-factor} and
\eqref{eq:outgoing-chain-total-length} give
\begin{equation*}
\prod_{k=1}^n C(P_k)
=
e^{-\varepsilon u}
\prod_{k=1}^n C^\varepsilon(P_k)
\leq
e^{-\varepsilon T}
\prod_{k=1}^n C^\varepsilon(P_k).
\end{equation*}
Applying the comparison above to the remaining bulk patches and to the end
contribution gives
\begin{equation*}
W_t^\mu\ind\bb{L_{T,t}^c}
\leq
e^{-\varepsilon T}W_t^{\varepsilon,\mu}.
\end{equation*}
Taking expectations and using~\eqref{eq:death-removed-weight-comparison} proves
\eqref{eq:finite-late-interaction-bound}.

For the infinite-horizon estimate, the same patchwise comparison gives
$0\leq W\leq W^\varepsilon$. On~$\cb{\abs{\mathcal P}<\infty}$ the product
defining~$W^\varepsilon$ has only finitely many factors, so
$W_s^{\varepsilon,\mu_{\mb 1}}\to W^\varepsilon$ as~$s\to\infty$, while
$W^\varepsilon=0$ on the complementary event. Hence
\begin{equation*}
W^\varepsilon
\leq
\liminf_{s\to\infty}W_s^{\varepsilon,\mu_{\mb 1}},
\end{equation*}
and Fatou's lemma together with~\eqref{eq:death-removed-weight-comparison}
gives
\begin{equation}
\label{eq:limiting-death-removed-mass}
\E_A\Cb{W^\varepsilon}\leq1.
\end{equation}
On~$\cb{\abs{\mathcal P}<\infty}\cap L_T^c$, apply the ancestry comparison to
the first successful interaction after~$T$ to obtain
\begin{equation*}
W\ind\bb{L_T^c}
\leq
e^{-\varepsilon T}W^\varepsilon.
\end{equation*}
On~$\cb{\abs{\mathcal P}=\infty}$ the left side is zero. Taking
expectations and using~\eqref{eq:limiting-death-removed-mass} proves
\eqref{eq:limiting-late-interaction-bound}.
\end{proof}

\begin{proof}[Proof of Lemma~\ref{lem:no-late-relaxation}]
On~$L_{T,t}$ there is no new patch boundary after time~$T$. Thus
\begin{equation*}
\mathcal B_t=\mathcal B_T,
\qquad
\mathcal E_t=\cb{P^{\uparrow t}:P\in\mathcal E_T},
\end{equation*}
so the bulk patches are fixed at time~$T$ and each end patch simply extends to
time~$t$.

\begin{figure}[htbp]
\centering
% Schematic of the no-late-interaction event L_{T,t}.
% The visual conventions are inherited from Figure 1 (patch-geometry.tex):
% red striped patch boundaries, gray outgoing-terminal patches, thin black
% finite-patch boundaries, and black successful interactions.
\begingroup
\definecolor{patchfigred}{RGB}{175,32,38}
\definecolor{patchfigactive}{gray}{0.86}
\definecolor{patchfigoutline}{gray}{0.15}
\begin{tikzpicture}[x=0.63cm,y=0.66cm]
\path[use as bounding box] (-0.08,-0.15) rectangle (20.58,10.66);
\tikzset{
  patchfig site/.style={draw=patchfigoutline,line width=0.48pt},
  patchfig success/.style={draw=black,line width=0.86pt,
    -{Stealth[length=2.0mm,width=1.35mm]}},
  patchfig active/.style={fill=patchfigactive,draw=none},
  patchfig horizon/.style={draw=patchfigoutline,line width=0.72pt},
  patchfig boundary/.style={draw=patchfigoutline,line width=0.40pt},
  patchfig cut/.style={draw=patchfigoutline,line width=1.25pt,
    dash pattern=on 4.2pt off 2.6pt,line cap=round},
  patchfig small/.style={font=\scriptsize,text=patchfigoutline,inner sep=0pt}
}

% The striped vertical boundary used for every patch in Figure 1.  Drawing one
% continuous band along a site column is equivalent here because consecutive
% patches share the same vertical sides.
\newcommand{\ConvEndPatchBarrier}[3]{%
  \begin{scope}
    \clip (#1,#2) rectangle ({#1+1},#3);
    \foreach \k in {-1,...,68}{%
      \pgfmathsetmacro{\yy}{-0.045+0.165*\k}
      \fill[patchfigred]
        (#1,\yy)--(#1,{\yy+0.075})--({#1+0.105},{\yy+0.155})
        --({#1+0.105},{\yy+0.080})--cycle;
      \fill[patchfigred]
        ({#1+1},\yy)--({#1+1},{\yy+0.075})
        --({#1+0.895},{\yy+0.155})--({#1+0.895},{\yy+0.080})--cycle;}
  \end{scope}}

% A successful interaction cuts the source column and the two target columns.
\newcommand{\ConvEndInteractionBoundary}[2]{%
  \draw[patchfig boundary] ({#1-1},#2)--({#1+2},#2);}
\newcommand{\ConvEndInteractionMark}[2]{%
  \draw[patchfig success] ({#1+0.5},#2)--({#1-0.5},#2);
  \draw[patchfig success] ({#1+0.5},#2)--({#1+1.5},#2);
  \fill[black] ({#1+0.5},#2) circle[radius=1.25pt];}

% A patch ending at the source of a successful interaction has terminal type
% O.  As in Figure 1, these deterministically active patches are filled gray.
\foreach \x/\a/\b in {
  9/0.00/0.48,10/0.48/0.92,11/0.92/1.20,8/0.48/1.38,
  9/1.38/1.58,7/1.38/1.74,12/1.20/2.18,10/1.58/2.40,
  6/1.74/2.72,8/1.74/2.88,13/2.18/2.98,5/2.72/3.22,
  11/2.40/3.44,14/2.98/3.70,6/3.22/3.94,15/3.70/4.02,
  4/3.22/4.12,12/3.44/4.36,9/2.88/4.58,3/4.12/4.68,
  16/4.02/4.86,2/4.68/5.02,14/4.02/5.18,4/4.68/5.28,
  12/4.36/5.32,17/4.86/5.38,7/3.94/5.50,10/4.58/5.66,
  1/5.02/5.74,18/5.38/5.86,15/5.18/5.98,5/5.28/6.10
}{\path[patchfig active] (\x,\a) rectangle ({\x+1},\b);}

% Patch columns begin when their sites first enter the interaction cone.  The
% unequal starting times give the cone an irregular, asymmetric outline.
\foreach \x/\s in {
  0/5.74,1/5.02,2/4.68,3/4.12,4/3.22,
  5/2.72,6/1.74,7/1.38,8/0.48,9/0.00,
  10/0.00,11/0.92,12/1.20,13/2.18,14/2.98,
  15/3.70,16/4.02,17/4.86,18/5.38,19/5.86
}{\ConvEndPatchBarrier{\x}{\s}{10.50}}

% A site boundary begins when at least one of its neighboring site columns
% enters the cone.  No vertical grid is drawn where neither side is a patch.
\foreach \x/\s in {
  0/5.74,1/5.02,2/4.68,3/4.12,4/3.22,
  5/2.72,6/1.74,7/1.38,8/0.48,9/0.00,
  10/0.00,11/0.00,12/0.92,13/1.20,14/2.18,
  15/2.98,16/3.70,17/4.02,18/4.86,19/5.38,
  20/5.86
}{\draw[patchfig site](\x,\s)--(\x,10.50);}
\draw[patchfig horizon](0,0)--(20,0);
\draw[patchfig horizon](0,10.50)--(20,10.50);

% Dense successful-interaction skeleton below T.  The outer interactions grow
% the cone one site at a time; the remaining interactions subdivide its bulk.
\foreach \s/\y in {
  9/0.48,10/0.92,11/1.20,8/1.38,9/1.58,7/1.74,
  12/2.18,10/2.40,6/2.72,8/2.88,13/2.98,5/3.22,
  11/3.44,14/3.70,6/3.94,15/4.02,4/4.12,12/4.36,
  9/4.58,3/4.68,16/4.86,2/5.02,14/5.18,4/5.28,
  12/5.32,17/5.38,7/5.50,10/5.66,1/5.74,18/5.86,
  15/5.98,5/6.10
}{\ConvEndInteractionBoundary{\s}{\y}}

% Draw interactions last so that the black arrows and source dots remain
% legible over the gray fills and striped patch boundaries.
\foreach \s/\y in {
  9/0.48,10/0.92,11/1.20,8/1.38,9/1.58,7/1.74,
  12/2.18,10/2.40,6/2.72,8/2.88,13/2.98,5/3.22,
  11/3.44,14/3.70,6/3.94,15/4.02,4/4.12,12/4.36,
  9/4.58,3/4.68,16/4.86,2/5.02,14/5.18,4/5.28,
  12/5.32,17/5.38,7/5.50,10/5.66,1/5.74,18/5.86,
  15/5.98,5/6.10
}{\ConvEndInteractionMark{\s}{\y}}

% T cuts through every end patch: it is a guide, not a patch boundary.  There
% are no interactions or horizontal subdivisions between T and t.
\draw[patchfig cut] (0,6.35)--(20,6.35);
\node[patchfig small,anchor=west] at (20.18,0) {$0$};
\node[patchfig small,anchor=west] at (20.18,6.35) {$T$};
\node[patchfig small,anchor=west] at (20.18,10.50) {$t$};
\end{tikzpicture}
\endgroup

\caption{On~$L_{T,t}$, no successful interaction occurs in~$(T,t]$. The bulk
patches are unchanged after~$T$, while every end patch present at~$T$ extends
to time~$t$.}
\label{fig:no-late-end-patches}
\end{figure}
\FloatBarrier

Fix the successful-interaction skeleton up to time~$T$. For
$P\in\mathcal E_T$, let~$L_{T,t}^P$ be the event that the one-site evolution
continuing~$P$ produces no successful interaction in~$(T,t]$. In Appendix~\ref{app:no-late-continuation} we prove that these events are conditionally independent, so
\begin{equation*}
L_{T,t}=\bigcap_{P\in\mathcal E_T}L_{T,t}^P,
\qquad
\CP{L_{T,t}}{\mathcal G_T}
=
\prod_{P\in\mathcal E_T}
\CP{L_{T,t}^P}{\mathcal G_T}.
\end{equation*}
Additionally, for~$P$ based at~$i$ we have
\begin{equation}
\label{eq:main-local-end-relaxation}
\CP{L_{T,t}^P}{\mathcal G_T}
C\bb{z,P^{\uparrow t}}
=
C\bb{\psi_i(t-T,z),P},
\end{equation}
where
\begin{equation*}
\psi_i(t-T,z)
=
\rho_i(\mb 0)
+
e^{-\lambda_i(\mb 0)(t-T)}\bb{z-\rho_i(\mb 0)}.
\end{equation*}
The pure-death assumption gives~$\lambda_i(\mb 0)\geq\varepsilon$. Since
$C(z,P)$ is affine in~$z$, the conditional end-patch contribution in
\eqref{eq:main-local-end-relaxation} converges to
$C\bb{\rho_i(\mb 0),P}$ at rate at least~$e^{-\varepsilon(t-T)}$.

Write
\begin{equation*}
B_T=\prod_{P\in\mathcal B_T}C(P),
\qquad
B_T^\varepsilon=\prod_{P\in\mathcal B_T}C^\varepsilon(P).
\end{equation*}
The conditional finite- and infinite-horizon contributions are therefore
\begin{align}
\label{eq:no-late-conditional-weights}
\CE{W_t^\mu\ind\bb{L_{T,t}}}{\mathcal G_T}
&=
B_T
\mu\bb{
 \prod_{P\in\mathcal E_T}\CP{L_{T,t}^P}{\mathcal G_T}
 C\bb{\eta\bb{i(P)},P^{\uparrow t}}
},\notag
\\
\CE{W\ind\bb{L_T}}{\mathcal G_T}
&=
B_T
\prod_{P\in\mathcal E_T}
C\bb{\rho_{i(P)}(\mb 0),P}.
\end{align}
Each factor under~$\mu$ in~\ref{eq:no-late-conditional-weights} converges
exponentially to the corresponding factor in the limiting product. To bound the error we use a telescoping estimate. The calculation details are in the Appendix, in the proof of 
Lemma~\ref{lem:end-factor-relaxation}, where we show
\begin{align*}
&\abs{
\CE{W_t^\mu\ind\bb{L_{T,t}}}{\mathcal G_T}
-
\CE{W\ind\bb{L_T}}{\mathcal G_T}
}
\\
&\qquad\leq
e^{-\varepsilon (t-T)}B_T
\CP{L_{T,t}}{\mathcal G_T}
\mu\bb{
 \prod_{P\in\mathcal E_T}
 C^\varepsilon\bb{\eta\bb{i(P)},P^{\uparrow t}}
}
\\
&\qquad\quad+
e^{-\varepsilon (t-T)}B_T\abs{\mathcal E_T}
\prod_{P\in\mathcal E_T}
C\bb{\rho_{i(P)}(\mb 0),P}.
\end{align*}
The first term is controlled by the finite-horizon contribution for the system
with pure deaths removed. The second accumulates one error for each end patch.
The event~$E_T^R$ is~$\mathcal G_T$-measurable, so by the tower property
\begin{align*}
&\abs{
 \E_A\Cb{W_t^\mu\ind\bb{E_T^R\cap L_{T,t}}}
 -
 \E_A\Cb{W\ind\bb{E_T^R\cap L_T}}
}
\\
&\quad\leq
e^{-\varepsilon (t-T)}
\E_A\Cb{
 \ind\bb{E_T^R}B_T
 \CP{L_{T,t}}{\mathcal G_T}
 \mu\bb{
  \prod_{P\in\mathcal E_T}
  C^\varepsilon\bb{\eta\bb{i(P)},P^{\uparrow t}}
 }
}
\\
&\qquad+
e^{-\varepsilon (t-T)}
\E_A\Cb{
 \ind\bb{E_T^R}B_T\abs{\mathcal E_T}
 \prod_{P\in\mathcal E_T}
 C\bb{\rho_{i(P)}(\mb 0),P}
}.
\end{align*}
For the first term,~$B_T\leq B_T^\varepsilon$ and the same formula
as~\eqref{eq:no-late-conditional-weights} for the system with pure deaths
removed give
\begin{equation*}
\E_A\Cb{
 \ind\bb{E_T^R}B_T
 \CP{L_{T,t}}{\mathcal G_T}
 \mu\bb{
  \prod_{P\in\mathcal E_T}
  C^\varepsilon\bb{\eta\bb{i(P)},P^{\uparrow t}}
 }
}
\leq
\E_A\Cb{
 W_t^{\varepsilon,\mu}\ind\bb{E_T^R\cap L_{T,t}}
}
\leq1.
\end{equation*}
For the second term,~$\abs{\mathcal E_T}\leq\abs R$ on
$E_T^R$, so~\eqref{eq:no-late-conditional-weights} gives
\begin{equation*}
\E_A\Cb{
 \ind\bb{E_T^R}B_T\abs{\mathcal E_T}
 \prod_{P\in\mathcal E_T}
 C\bb{\rho_{i(P)}(\mb 0),P}
}
\leq
\abs R\,
\E_A\Cb{
 W\ind\bb{E_T^R\cap L_T}
}
\leq\abs R.
\end{equation*}
Here the last inequalities use~\eqref{eq:death-removed-weight-comparison} and
$0\leq W\leq W^\varepsilon$ together with
\eqref{eq:limiting-death-removed-mass}. Combining the two bounds proves
\eqref{eq:no-late-relaxation-bound}.
\end{proof}

\begin{proof}[Proof of Lemma~\ref{lem:spatial-confinement}]
On~$E_T^R$, every successful interaction up to time~$T$ stays inside~$R$ so, if~$A\subseteq R$, the dual process remains inside~$R$ up to time~$T$. Fixing the spins outside~$R$ to~$0$ does not
change the contribution of such a realization. After~$T$ the dual process evolves without constraint. Proposition~\ref{prop:confined-interaction-identity} makes this argument precise by undoing the duality under the~$E_T^R$ event. We obtain
\begin{equation*}
\E_A\Cb{W_t^\mu\ind\bb{E_T^R}}
=
(\mu P_{t-T}P_T^{R,0})(\chi_A).
\end{equation*}

\begin{figure}[htbp]
\centering
% Schematic of the spatial-escape event (E_T^R)^c.
% The panel is sized to sit beside two further convergence-event panels.
\begingroup
\definecolor{conveventred}{RGB}{175,32,38}
\definecolor{conveventoutline}{gray}{0.18}
\definecolor{conveventcone}{gray}{0.82}
\definecolor{conveventfuture}{gray}{0.93}
\definecolor{conveventregion}{gray}{0.975}
\definecolor{conveventsecondary}{gray}{0.48}
\begin{tikzpicture}[x=0.54cm,y=0.54cm]
\path[use as bounding box] (0.18,-1.02) rectangle (7.12,7.34);
\tikzset{
  convevent boundary/.style={draw=conveventoutline,line width=0.72pt,line join=round},
  convevent guide/.style={draw=conveventsecondary,line width=0.52pt,dash pattern=on 2.3pt off 2.1pt},
  convevent baseline/.style={draw=conveventoutline,line width=0.52pt},
  convevent axis/.style={draw=conveventoutline,line width=0.62pt,-{Stealth[length=1.8mm,width=1.2mm]}},
  convevent tick/.style={draw=conveventoutline,line width=0.52pt},
  convevent label/.style={font=\scriptsize,text=conveventoutline,inner sep=1pt},
  convevent secondary label/.style={font=\scriptsize,text=conveventsecondary,inner sep=1pt}
}

% The spacetime cylinder R x [0,T].
\path[fill=conveventregion] (1.90,0) rectangle (4.90,4.60);

% The part of the cone after T is shown lightly because its shape is
% irrelevant to the event (E_T^R)^c.
\path[fill=conveventfuture]
  (0.98,4.60)
  .. controls (0.83,5.30) and (0.68,6.24) .. (0.54,7.00)
  -- (5.54,7.00)
  .. controls (5.18,6.22) and (4.76,5.28) .. (4.46,4.60)
  -- cycle;

% Before T the cone grows normally on the right.  On the left it has a
% short, atypically rapid expansion and leaves R.
\path[fill=conveventcone]
  (3.18,0)
  .. controls (3.10,0.94) and (3.02,1.86) .. (2.92,2.70)
  .. controls (2.86,3.20) and (2.78,3.57) .. (2.58,3.82)
  .. controls (2.39,4.06) and (1.34,4.14) .. (0.98,4.60)
  -- (4.46,4.60)
  .. controls (4.36,4.03) and (4.26,3.50) .. (4.17,2.92)
  .. controls (4.04,2.04) and (3.91,0.93) .. (3.80,0)
  -- cycle;

% Highlight exactly the pre-T part of the cone outside R.  Clipping to the
% cone keeps the red and black boundaries identical.
\begin{scope}
  \clip
    (3.18,0)
    .. controls (3.10,0.94) and (3.02,1.86) .. (2.92,2.70)
    .. controls (2.86,3.20) and (2.78,3.57) .. (2.58,3.82)
    .. controls (2.39,4.06) and (1.34,4.14) .. (0.98,4.60)
    -- (4.46,4.60)
    .. controls (4.36,4.03) and (4.26,3.50) .. (4.17,2.92)
    .. controls (4.04,2.04) and (3.91,0.93) .. (3.80,0)
    -- cycle;
  \path[fill=conveventred,fill opacity=0.42]
    (0.20,0) rectangle (1.90,4.60);
\end{scope}

% Continuous outer boundary of the cone.
\draw[convevent boundary]
  (3.18,0)
  .. controls (3.10,0.94) and (3.02,1.86) .. (2.92,2.70)
  .. controls (2.86,3.20) and (2.78,3.57) .. (2.58,3.82)
  .. controls (2.39,4.06) and (1.34,4.14) .. (0.98,4.60)
  .. controls (0.83,5.30) and (0.68,6.24) .. (0.54,7.00)
  -- (5.54,7.00)
  .. controls (5.18,6.22) and (4.76,5.28) .. (4.46,4.60)
  .. controls (4.36,4.03) and (4.26,3.50) .. (4.17,2.92)
  .. controls (4.04,2.04) and (3.91,0.93) .. (3.80,0)
  -- cycle;

% Guides for R and the cut time T.
\draw[convevent baseline] (0.30,0)--(5.82,0);
\draw[convevent guide] (1.90,0)--(1.90,4.60);
\draw[convevent guide] (4.90,0)--(4.90,4.60);
\draw[convevent guide] (0.30,4.60)--(6.34,4.60);
\node[convevent secondary label,anchor=south west] at (2.02,0.12) {$R$};

% Time axis and panel label.
\draw[convevent axis] (6.34,0)--(6.34,7.24);
\draw[convevent tick] (6.24,0)--(6.44,0);
\draw[convevent tick] (6.24,4.60)--(6.44,4.60);
\draw[convevent tick] (6.24,7.00)--(6.44,7.00);
\node[convevent label,anchor=west] at (6.50,0) {$0$};
\node[convevent label,anchor=west] at (6.50,4.60) {$T$};
\node[convevent label,anchor=west] at (6.50,7.00) {$t$};
\node[convevent label,anchor=north] at (3.36,-0.24) {$(E_T^R)^c$};
\end{tikzpicture}
\endgroup

\caption{The event~$(E_T^R)^c$: before time~$T$, the interaction cone leaves
$R$ due to unlikely growth spurt. The red region marks the part of the cone outside~$R$.}
\label{fig:convergence-spatial-escape}
\end{figure}
\FloatBarrier

The patch representation gives
\begin{equation*}
\E_A\Cb{W_t^\mu}
=
(\mu P_t)(\chi_A)
=
(\mu P_{t-T}P_T)(\chi_A).
\end{equation*}
Therefore
\begin{align*}
0
&\leq
\E_A\Cb{W_t^\mu\ind\bb{(E_T^R)^c}}
\\
&=
(\mu P_{t-T})\bb{\bb{P_T-P_T^{R,0}}\chi_A}
\\
&\leq
\norm{\bb{P_T-P_T^{R,0}}\chi_A}_\infty\\
&=
\rho_A(T,R),
\end{align*}
which proves~\eqref{eq:finite-horizon-confinement-bound}.

On~$\cb{\abs{\mathcal P}<\infty}$,
$W_u^{\mu}\to W$ as~$u\to\infty$, and on~$\cb{\abs{\mathcal P}=\infty}$ we have~$W=0$, so Fatou's lemma gives
\begin{equation*}
0
\leq
\E_A\Cb{W\ind\bb{(E_T^R)^c}}
\leq
\liminf_{u\to\infty}\E_A\Cb{W_u^\mu\ind\bb{(E_T^R)^c}}
\leq
\rho_A(T,R),
\end{equation*}
which is~\eqref{eq:limiting-confinement-bound}.
\end{proof}

\begingroup
\color{red}
\subsection{Ergodicity under orientability}
\label{subsec:oriented-ergodicity}


% ============================================================================
% BLURB A. What remains after Theorem C?
%
% Reader question:
%   Why should approaching M_* imply convergence from an arbitrary state?
% Role in the proof:
%   State the reduction before introducing restricted dynamics or conditioning.
% Drafting target:
%   One opening paragraph. Explain local closeness; do not suggest that a
%   single finite burn-in time puts the full law in M_* or M_{-,1}.
% ============================================================================
The statement of Corollary~\ref{cor:oriented-ergodicity} is that orientability suffices to strengthen 
Theorem~\ref{thm:common-invariant-limit} to ergodicity. The main idea is as follows: since all measures in~$\mathcal{M}_*$ converge to~$\pi$ then ,to prove ergodicity, it suffices to show that for any configuration~$\eta$ its evolution~$\delta_\eta P_t$ converges towards~$\mathcal{M}_*$. More precisely, by compactness and the Feller property, it
suffices to show that~$\delta_\eta P_t$ approaches~$\mathcal M_*$ on finite
sets, in the sense that
\begin{equation*}
\liminf_{t\to\infty}\inf_\eta P_t(\chi_A^*)(\eta)\geq0
\qquad\text{for every finite }A\subseteq\Lambda.
\end{equation*}
To obtain this bound, choose a finite set~$I=\bar A$ as in the definition
of orientability and consider the dynamics of the spin system on~$I$, conditional on a particular realization of the trajectory on the (autonomous) set~$U(A)$. This is the original spin system, but now on a finite set and with time-inhomogenous boundary conditions. This restricted system is still patch positive and, with the original
centering profile~$p^*$, preserves the class~$M_*^I$ of probability measures on~$\cb{0,1}^I$ with non-negative centered moments. The boundary conditions are not homogenous in time, so the restricted system does not converge to any invariant measure, but it does converge towards the class~$\mathcal{M}_*^I$. Finiteness of~$I$ and presence of noise are crucial here, because they guarantee that a positive part of the measure enters~$\mathcal{M}_*^I$ after running the restricted process for any positive time. This proves the convergence towards~$\mathcal{M}_*$, in the above sense, and hence ergodicity.
Given uniform orientability, we make convergence rate quantitative and show it is exponential.

\begin{proof}[Proof of Corollary~\ref{cor:oriented-ergodicity}]
    


% ============================================================================
% BLURB B. What process do we obtain after fixing the exterior?
%
% Reader question:
%   What is the restricted process, and which class of measures does it preserve?
% Role in the proof:
%   Define the finite system and M_*^I at the ORIGINAL profile p^star.
% Drafting target:
%   Two short paragraphs. State preservation with a reference to the appendix
%   lemma; keep its algebraic proof there. No conditioning is used here.
% ============================================================================
Fix a nonempty finite set~$A$ and take~$I=\bar A$ and~$U=U(A)$ from
orientability. For~$\xi\in\cb{0,1}^U$, let~$\eta^{U,\xi}$ agree with~$\xi$
on~$U$ and with~$\eta$ elsewhere. Since~$I\cup U$ is autonomous, fixing
the spins on~$U$ defines a spin system on~$\cb{0,1}^I$, with generator
and rates
\begin{align*}
\Lcal^{I,\xi}f(\eta)&=\Lcal f(\eta^{U,\xi}),\\
\lambda_i^{I,\xi}(\eta)&=\lambda_i(\eta^{U,\xi}),
\qquad
\rho_i^{I,\xi}(\eta)=\rho_i(\eta^{U,\xi})
\quad(i\in I).
\end{align*}
Denote the restricted semigroup by~$P_t^{I,\xi}$.

With the original centering profile~$(p_i^\star)_{i\in I}$, define
\begin{equation*}
\mathcal M_*^I
=\cb{\mu\in\mathcal P\bb{\cb{0,1}^I}:
       \mu(\chi_B^*)\geq0\text{ for every }B\subseteq I}.
\end{equation*}
Lemma~\ref{lem:boundary-moment-preservation} in
Appendix~\ref{app:boundary-moment-preservation} gives
\begin{equation*}
\mathcal M_*^I P_t^{I,\xi}\subseteq\mathcal M_*^I.
\end{equation*}
This holds for every boundary condition~$\xi$. The restricted system is also patch positive, though we do not use this fact directly.

% ============================================================================
% BLURB C. Where does a positive part in M_*^I come from?
%
% Reader question:
%   What event erases the initial state on a finite set?
% Role in the proof:
%   Obtain a lower bound uniform in the initial state and prescribed boundary.
% Drafting target:
%   One paragraph describing the reset event, followed by the bound.
% Correction:
%   A coefficient of delta_e cannot generally tend to one. At small times,
%   forcing every site costs order s^{|I|}; a bound 1-exp(-alpha_I s) for
%   every s>0 is generally false. Later extraction uses fixed unit intervals.
% ============================================================================
\paragraph{Blurb C. A positive part after one time interval.}
 Use~$\lambda_i\rho_i\geq\varepsilon p_i^\star$ and the pure-death
component to retain independent resets to~$e_i$ at rate~$r_i$, where
\begin{equation*}
e_i=\ind\bb{p_i^\star>0},
\qquad
r_i=
\begin{cases}
\varepsilon p_i^\star,&p_i^\star>0,\\
\varepsilon,&p_i^\star=0.
\end{cases}
\end{equation*}
Write~$\overline\lambda=\sup_{i,\eta}\lambda_i(\eta)$.
Describe the event that each retained reset clock rings and no remaining
update mark occurs in~$I$ during an interval of length~$s$. The required
bound, for every~$x\in\cb{0,1}^I$, is
\begin{equation*}
P_s^{I,\xi}(x,\cdot)\geq\alpha_I(s)\delta_e,
\qquad
\alpha_I(s)=e^{-\overline\lambda\abs I s}
\prod_{i\in I}(1-e^{-r_i s})>0
\quad(s>0),
\qquad
\delta_e\in\mathcal M_*^I.
\end{equation*}
The same event works if the prescribed boundary changes during the interval.
Retain the constant~$a_I=\alpha_I(1)>0$ for repeated extraction.
Distinguish this fixed positive part from the total mass in~$\mathcal M_*^I$,
which will approach~$1$ after previously extracted parts have evolved.

% ============================================================================
% BLURB D. How do boundary conditioning and repeated extraction fit together?
%
% Reader question:
%   Why does revealing the exterior give the prescribed-boundary process,
%   and why do arbitrarily frequent boundary changes cause no loss?
% Role in the proof:
%   Derive the conditional centered-moment lower bound with the correct
%   chronological order and geometric mixture weights.
% Drafting target:
%   Separate the conditioning paragraph from the extraction paragraph.
% ============================================================================
\paragraph{Blurb D. Conditioning and repeated extraction.}
Define
\begin{equation*}
\mathcal G_U=\sigma\cb{\eta_s(j):j\in U,\ s\geq0},
\qquad
\partial I=N_*(I)\setminus I\subseteq U.
\end{equation*}
For a deterministic initial configuration~$\eta$, justify conditioning on
$\mathcal G_U$ using autonomy of~$U$ and~$I\cup U$. Only the finitely many
spins in~$\partial I$ affect the restricted dynamics. If their successive
values on~$[0,t]$ are~$\xi_1,\ldots,\xi_n$ for durations
$t_1,\ldots,t_n$ with~$\sum_k t_k=t$, write
\begin{equation*}
\E_\eta\Cb{\chi_A^*(\eta_t)\mid\mathcal G_U}
=
\bb{P_{t_1}^{I,\xi_1}\cdots P_{t_n}^{I,\xi_n}\chi_A^*}(\eta|_I).
\end{equation*}
Here each~$\xi_k$ may be extended arbitrarily to the rest of~$U$.
Denote the resulting evolution along the revealed boundary path~$b$
by~$Q^b_{s,t}$. Lemma~\ref{lem:boundary-moment-preservation} gives
preservation of~$\mathcal M_*^I$ under this finite composition.

Perform extraction on fixed unit intervals, independently of the boundary
change times. Blurb~C gives Markov kernels~$R_k^b$ such that
\begin{equation*}
Q^b_{k-1,k}(x,\cdot)=a_I\delta_e+(1-a_I)R_k^b(x,\cdot).
\end{equation*}
For~$m=\lfloor t\rfloor$, expand according to the first extracted part:
\begin{align*}
\delta_{\eta|_I}Q^b_{0,t}
&=\sum_{k=1}^m a_I(1-a_I)^{k-1}\delta_e Q^b_{k,t}
 +(1-a_I)^m\nu_t^b,\\
\nu_t^b
&=\delta_{\eta|_I}R_1^b\cdots R_m^b Q^b_{m,t}.
\end{align*}
The measure~$\nu_t^b$ is a probability measure, and every term
$\delta_e Q^b_{k,t}$ belongs to~$\mathcal M_*^I$. The target is therefore
\begin{equation*}
\E_\eta\Cb{\chi_B^*(\eta_t)\mid\mathcal G_U}
\geq-\norm{\chi_B^*}_\infty(1-a_I)^{\lfloor t\rfloor}
\qquad(B\subseteq I).
\end{equation*}
Average over~$\mathcal G_U$ and take~$B=A$.

% ============================================================================
% BLURB E. Why does the moment bound imply ergodicity?
%
% Reader question:
%   How does a statement for each finite set become convergence of the law?
% Role in the proof:
%   Supply the limiting argument; local cone approximation alone must not be
%   treated as an application of Theorem C without this step.
% Drafting target:
%   One paragraph for compactness/Feller and Theorem C. Keep the finite-volume
%   M_{-,1} observation separate, and do not require it for this argument.
% ============================================================================
\paragraph{Blurb E. From the burn-in estimate to ergodicity.}
Explain that every weak subsequential limit of
$\delta_{\eta_n}P_{t_n}$ with~$t_n\to\infty$ belongs to~$\mathcal M_*$:
apply Blurb~D to each finite~$A$. For fixed~$s$, apply this observation to
$\delta_{\eta_n}P_{t_n-s}$ and use compactness and the Feller property to
justify the first inequality below:
\begin{equation*}
\limsup_{t\to\infty}\sup_\eta\abs{P_tf(\eta)-\pi(f)}
\leq
\sup_{\mu\in\mathcal M_*}\abs{(\mu P_s)(f)-\pi(f)}
\leq K_f(1+s)^D e^{-\varepsilon s/2}
\longrightarrow0
\quad(s\to\infty).
\end{equation*}
Conclude convergence from every initial law and uniqueness of~$\pi$.

Record separately the answer to the finite-time entry question.
For fixed~$I$, uniform deaths imply~$p_i^\star<1$, so
\begin{equation*}
(1-a_I)^{\lfloor t\rfloor}
\leq
\min_{\vn\neq B\subseteq I}
\frac{\prod_{i\in B}(1-p_i^\star)}{\norm{\chi_B^*}_\infty}
\quad\Longrightarrow\quad
(\delta_\eta P_t)|_I\in\mathcal M_{-,1}^I,
\end{equation*}
where~$\mathcal M_{-,1}^I=\cb{\mu:(\mu+\delta_{\mb1_I})/2\in\mathcal M_*^I}$.
The minimum is positive. This gives a time depending on~$I$, uniformly
in~$\eta$; it does not establish entry into the global~$\mathcal M_{-,1}$.

% ============================================================================
% BLURB F. What quantitative information does uniform orientability add?
%
% Reader question:
%   Why is there now a common exponential rate for every local observable?
% Role in the proof:
%   Replace finite-set extraction by a boundary-uniform coupling estimate,
%   then combine burn-in with finite propagation and Theorem C.
% Drafting target:
%   Explain the choice of region before the coupling, and the coupling
%   before the final estimate. Leave the detailed path count to drafting.
% ============================================================================
\paragraph{Blurb F. Uniform orientability and an exponential rate.}
For~$f$ supported on~$A$, choose a finite-propagation speed~$v$ for error
$e^{-\varepsilon s/2}$, independent of~$A$. With the height function~$h$
and increment bound~$r$, use
\begin{equation*}
H_s=\max_{i\in A}h(i)+r\lceil vs\rceil,
\qquad
I_s=\bigcup_{n\geq0}B(A,n)\cap\cb{i:h(i)\leq H_s},
\qquad
U_s=\cb{i:h(i)>H_s}.
\end{equation*}
Explain autonomy and the geometric bounds
\begin{equation*}
B(A,\lceil vs\rceil)\subseteq I_s,
\qquad
\abs{I_s}\leq C_A(1+s)^D,
\qquad
\ell_s\leq d_A+rv s,
\end{equation*}
where~$\ell_s$ is the maximum number of vertices in a directed path in~$I_s$.
Couple the restrictions from~$x$ and~$\mb1_{I_s}$ with the same boundary,
pure-death clocks, and remaining graphical randomness. Explain how agreement
passes from greater to smaller heights. With
$d=\max\cb{1,\sup_i\abs{N(i)}}$, the estimate to prove is
\begin{equation*}
\sup_{x,b}d_{\mathrm{TV}}\bb{\delta_xQ^b_{0,u},
\delta_{\mb1_{I_s}}Q^b_{0,u}}
\leq2\abs{I_s}(2d)^{\ell_s}e^{-\varepsilon u/2}.
\end{equation*}
Here~$d_{\mathrm{TV}}$ is the supremum over events. Explain why averaging
the second law over the boundary and extending by all~$1$'s gives a law
in~$\mathcal M_*$. Compare~$P_sf$ with~$P_s^{I_s,0}f$ using
Lemma~\ref{lem:finite-propagation}, then apply
Theorem~\ref{thm:common-invariant-limit} to obtain
\begin{equation*}
\sup_\eta\abs{P_{u+s}f(\eta)-\pi(f)}
\leq C_f(1+s)^D
\bb{e^{-\varepsilon s/2}+e^{rv\log(2d)s-\varepsilon u/2}}.
\end{equation*}
Finish with~$u=\bb{1+2rv\log(2d)/\varepsilon}s$ and explain why the
resulting exponential rate is independent of~$f$.
\par
\end{proof}
\endgroup
